\expandafter\ifx\csname CAPETINCLUDEMODE\endcsname\relax
\documentclass[11pt,a4paper]{article}
% preamble.tex — packages et commandes communes a toutes les fiches CAPET SII IE
\input{glyphtounicode}
\pdfgentounicode=1
\usepackage{cmap}
\usepackage[T1]{fontenc}
\usepackage[utf8]{inputenc}
\usepackage[french]{babel}
% babel[french] active a \begin{document} des raccourcis de ponctuation (:;!?)
% qui cassent la syntaxe to[...] de circuitikz (erreur "+ or - expected") ;
% on les desactive, ce qui ne change que l'espacement fin avant ces signes.
\AtBeginDocument{\shorthandoff{:;!?}}
\usepackage[a4paper, top=1.6cm, bottom=1.6cm, left=1.4cm, right=1.4cm]{geometry}
\setlength{\headheight}{14.5pt}
\usepackage{amsmath, amssymb, amsfonts}
\usepackage{siunitx}
\sisetup{locale=FR, output-decimal-marker={,}, inter-unit-product=\ensuremath{{}\cdot{}},
         group-minimum-digits=4}
% Unites non fournies par siunitx, declarees ici une fois pour toutes (fiches + exercices).
% Les gardes \ifdefined evitent tout conflit avec une redeclaration locale restee dans un
% fichier : ne pas retirer les gardes, ne pas redeclarer ces unites dans les fichiers.
\DeclareSIUnit{\jour}{j}
\ifdefined\var\else\DeclareSIUnit{\var}{var}\fi   % puissance reactive
\ifdefined\vad\else\DeclareSIUnit{\vad}{VAD}\fi   % puissance deformante
\ifdefined\tr\else\DeclareSIUnit{\tr}{tr}\fi      % tour (tr/min)
\usepackage{array, booktabs, multirow, tabularx}
\usepackage{xcolor}
\usepackage[most, breakable]{tcolorbox}
\usepackage{multicol}
\usepackage{enumitem}
\usepackage{mathtools}
\usepackage{lmodern}
% Refonte visuelle 2026 (direction "moderne coloré") : Lato pour le texte,
% maths sans-serif assorties via newtxsf (charge apres amsmath/amssymb).
\usepackage[default]{lato}
\usepackage{newtxsf}
\usepackage{microtype}
\usepackage{fancyhdr}
\usepackage{tikz}
\usetikzlibrary{positioning,arrows.meta}
\usepackage{circuitikz}
\usepackage{lastpage}

% ── Couleurs par domaine (palette modernisée 2026) ────────────────────
\definecolor{colElec}{RGB}{29, 78, 216}        % Électrotechnique — bleu vif
\definecolor{colElecLight}{RGB}{219, 234, 254} % fond bleu clair
\definecolor{colAuto}{RGB}{21, 128, 61}        % Automatique — vert
\definecolor{colMeca}{RGB}{194, 65, 12}        % Mécanique — orange
\definecolor{colTherm}{RGB}{185, 28, 28}       % Thermique — rouge
\definecolor{colQT}{RGB}{232, 244, 233}        % (hérité, plus utilisé par questiontype)
\definecolor{colJury}{RGB}{255, 248, 225}      % fond jury (crème)
\definecolor{colJuryFrame}{RGB}{202, 138, 4}   % bordure/bandeau jury (ambre)
\definecolor{colGrayLight}{RGB}{245, 245, 247} % fond gris neutre
\definecolor{colSeparator}{RGB}{180, 180, 180} % gris séparateur
\definecolor{colExo}{RGB}{8, 145, 178}          % Banque d'exercices — cyan (distinct des fiches)
% ── Encre et gris de la refonte ───────────────────────────────────────
\definecolor{colInk}{RGB}{23, 32, 42}          % texte courant (quasi-noir)
\definecolor{colSlate}{RGB}{83, 98, 118}       % gris ardoise (méta, en-têtes)
\definecolor{colRule}{RGB}{203, 213, 225}      % filets clairs
\definecolor{colJuryBg}{RGB}{255, 248, 225}    % fond crème jury
\definecolor{colPresence}{RGB}{83, 98, 118}    % bloc 5 (présence dans les sujets)
\AtBeginDocument{\color{colInk}}

% ── Style encadré exercice + corrigé (banque d'exercices) ─────────────
\tcbset{
  exoBox/.style={
    breakable,
    enhanced,
    colback=colExo!8!white,
    colframe=colExo,
    colbacktitle=colExo,
    coltitle=white,
    fonttitle=\bfseries\small,
    left=4pt, right=4pt, top=3pt, bottom=3pt,
    before skip=5pt, after skip=5pt,
  }
}

% ── Bascule enonce / corrige (banque d'exercices) ─────────────────────
% \CAPETCORRIGE est defini en ligne de commande (pdflatex "\def\CAPETCORRIGE{1}\input{...}")
% pour compiler la variante corriges d'un meme fichier source ; par defaut (non defini),
% la compilation produit la variante enonces. \CAPETSHOW{enonce}{corrige} choisit la branche.
%
% TROIS MODES (refonte 2026-09-20, manuel de revision) :
%   - aucun drapeau   -> imprime l'enonce            (variante enonces, INCHANGE)
%   - \CAPETCORRIGE   -> imprime le corrige          (variante corriges, INCHANGE)
%   - \CAPETMANUEL    -> imprime l'enonce ET met le corrige de cote pour l'annexe
%
% Le mode manuel stocke chaque corrige dans une macro (\csgdef), sans ecriture
% fichier ni re-tokenisation : le contenu (circuitikz, accents, %) traverse intact.
% L'annexe les rejoue ensuite dans l'ordre via \CAPETANNEXE.
\usepackage{etoolbox}

\newcounter{CAPETexo}          % numero continu d'exercice dans le manuel
\setcounter{CAPETexo}{0}
\newcounter{CAPETannexe}       % compteur de relecture en annexe

% Titre courant de la notion en cours (pose par \CAPETNOTION), sert d'intertitre annexe
\newcommand{\CAPETnotioncode}{}
\newcommand{\CAPETnotiontitre}{}

% \CAPETNOTION{code}{titre} : ouvre une notion dans le manuel (memorise le libelle,
% et marque le rang du premier exercice de la notion pour l'intertitre d'annexe).
\newcommand{\CAPETNOTION}[2]{%
  \renewcommand{\CAPETnotioncode}{#1}%
  \renewcommand{\CAPETnotiontitre}{#2}%
  \csgdef{CAPETnotionstart\the\numexpr\value{CAPETexo}+1\relax}{#1 -- #2}%
}

% ── \degree robuste ───────────────────────────────────────────────────
% Certains fichiers d'exercices utilisent \degree. Ni le preambule ni siunitx
% ne le definissent, et \providecommand local ne suffit pas : sous hyperref,
% \degree existe deja comme commande factice (\ERROR), donc \providecommand
% s'abstient et l'appel echoue a la compilation du manuel (erreur invisible
% en compilation isolee). On le (re)definit donc inconditionnellement ici,
% apres le chargement des packages, pour tout le corpus.
\AtBeginDocument{\providecommand{\degree}{}\renewcommand{\degree}{\ensuremath{^\circ}}}

% Libelle decoratif : les fichiers d'exercices utilisent aussi \CAPETSHOW dans
% leur \fancyhead et leur bloc de titre, pour afficher "(enonces)" ou "(corriges)".
% Ces appels-la ne sont PAS des exercices : en mode manuel ils ne doivent ni
% consommer un numero, ni alimenter l'annexe. On les reconnait a leur premier
% argument, strictement egal au libelle ci-dessous.
\newcommand{\CAPETLIBELLEENONCES}{(\'enonc\'es)}

\ifdefined\CAPETMANUEL
  % --- Mode manuel : enonce en place, corrige differe vers l'annexe ---
  \newcommand{\CAPETSHOW}[2]{%
    \def\CAPETargun{#1}%
    \ifx\CAPETargun\CAPETLIBELLEENONCES
      #1%            appel decoratif : rendu tel quel, hors numerotation
    \else
      \CAPETSHOWEXO{#1}{#2}%
    \fi
  }
  \newcommand{\CAPETSHOWEXO}[2]{%
    \stepcounter{CAPETexo}%
    \csgdef{CAPETcorrige\theCAPETexo}{#2}%
    \edef\CAPETtmpnum{\theCAPETexo}%
    #1%
    \par\vspace{2pt}%
    {\footnotesize\itshape\color{colExo}%
     \hyperref[capet:cor:\CAPETtmpnum]{Corrige de l'exercice \CAPETtmpnum\ $\rightarrow$ p.~\pageref{capet:cor:\CAPETtmpnum}}%
     \label{capet:eno:\CAPETtmpnum}}%
  }
  % \CAPETANNEXE : rejoue tous les corriges captures, dans l'ordre, avec intertitres
  \newcommand{\CAPETANNEXE}{%
    \setcounter{CAPETannexe}{0}%
    \loop\ifnum\value{CAPETannexe}<\value{CAPETexo}%
      \stepcounter{CAPETannexe}%
      \edef\CAPETcurnum{\theCAPETannexe}%
      \ifcsname CAPETnotionstart\CAPETcurnum\endcsname
        \par\medskip
        {\normalsize\bfseries\color{colExo}\csuse{CAPETnotionstart\CAPETcurnum}}%
        \par\nobreak\smallskip
      \fi
      \begin{tcolorbox}[exoBox, breakable,
        title={Corrige de l'exercice \CAPETcurnum\
               \hfill\normalfont\footnotesize
               \hyperref[capet:eno:\CAPETcurnum]{enonce p.~\pageref{capet:eno:\CAPETcurnum}}}]%
      \label{capet:cor:\CAPETcurnum}%
      \csuse{CAPETcorrige\CAPETcurnum}%
      \end{tcolorbox}%
    \repeat
  }
\else
  % --- Modes historiques, comportement strictement inchange ---
  \newcommand{\CAPETSHOW}[2]{\ifdefined\CAPETCORRIGE #2\else #1\fi}
  \newcommand{\CAPETANNEXE}{}
\fi

% ── Difficulté 1 à 3 (étoiles pleines, amssymb) ───────────────────────
\newcommand{\difficulte}[1]{%
  \ifcase#1\relax
  \or $\bigstar$\;niveau 1/3%
  \or $\bigstar\bigstar$\;niveau 2/3%
  \or $\bigstar\bigstar\bigstar$\;niveau 3/3%
  \fi
}

% ── Style encadré extrait de document technique (banque d'exercices) ──
\tcbset{
  docTech/.style={
    breakable,
    enhanced,
    colback=colGrayLight,
    colframe=colSeparator,
    fonttitle=\bfseries\small\itshape,
    coltitle=black,
    colbacktitle=colGrayLight,
    left=4pt, right=4pt, top=3pt, bottom=3pt,
    before skip=4pt, after skip=4pt,
  }
}

% ── Style tcolorbox de section principale ────────────────────────────
% Refonte 2026 : bandeau titre coloré arrondi + fond très pâle, sans cadre.
\tcbset{
  ficheSection/.style={
    breakable, enhanced jigsaw,
    frame hidden, boxrule=0pt, arc=1.8mm,
    colback=#1!4!white,
    colbacktitle=#1,
    coltitle=white,
    fonttitle=\sffamily\bfseries,
    left=7pt, right=7pt, top=5pt, bottom=5pt,
    toptitle=3pt, bottomtitle=3pt,
    before skip=10pt, after skip=8pt,
  }
}

% ── Style encadré question type ──────────────────────────────────────
% Carte blanche, bandeau sombre, liseré très clair.
\tcbset{
  questiontype/.style={
    breakable, enhanced jigsaw,
    colback=white,
    frame hidden, boxrule=0pt, arc=1.8mm,
    colbacktitle=colElec!78!black,
    coltitle=white,
    fonttitle=\sffamily\bfseries\small,
    left=7pt, right=7pt, top=4pt, bottom=4pt,
    toptitle=2.5pt, bottomtitle=2.5pt,
    before skip=8pt, after skip=8pt,
    borderline={0.6pt}{0pt}{colElec!12!white},
  }
}

% ── Style encadré jury ───────────────────────────────────────────────
% Bandeau ambre, fond crème, sans cadre.
\tcbset{
  jury/.style={
    breakable, enhanced jigsaw,
    colback=colJuryBg,
    frame hidden, boxrule=0pt, arc=1.8mm,
    colbacktitle=colJuryFrame,
    coltitle=white,
    fonttitle=\sffamily\bfseries,
    left=7pt, right=7pt, top=4pt, bottom=4pt,
    toptitle=3pt, bottomtitle=3pt,
    before skip=10pt, after skip=8pt,
  }
}

% ── En-tête / pied de page communs (refonte 2026) ─────────────────────
% #1 code fiche, #2 titre court, #3 couleur domaine
\newcommand{\ficheHeaders}[3]{%
  \fancyhead[L]{\small\sffamily{\bfseries\textcolor{#3}{#1}}\hspace{0.7em}\textcolor{colSlate}{#2}}%
  \fancyhead[R]{\small\sffamily\textcolor{colSlate}{CAPET SII IE\, --- \,Session 2026}}%
  \fancyfoot[C]{\small\sffamily\textcolor{colSlate}{\thepage\,/\,\pageref{LastPage}}}%
  \fancyhead[C]{}\fancyfoot[L]{}\fancyfoot[R]{}%
}

% ── Bloc en-tête de fiche (BLOC 0, refonte 2026) ──────────────────────
% #1 code  #2 titre  #3 couleur domaine  #4 domaine  #5 ligne présence  #6 sous-thèmes
\newcommand{\ficheEntete}[6]{%
  \begin{tcolorbox}[enhanced, breakable,
    frame hidden, boxrule=0pt, arc=2.5mm,
    interior style={left color=#3, right color=#3!75!black},
    left=9pt, right=9pt, top=7pt, bottom=7pt,
    before skip=0pt, after skip=6pt]
    {\color{white}\sffamily\bfseries\LARGE #1\quad #2}\\[4pt]
    {\color{#3!22!white}\sffamily\small \MakeUppercase{#4}\quad\textbar\quad #5}\\[1pt]
    {\color{#3!35!white}\sffamily\footnotesize Sous-thèmes : #6}
  \end{tcolorbox}%
}

% ── Mise en page ─────────────────────────────────────────────────────
\pagestyle{fancy}
\fancyhf{}
\renewcommand{\headrulewidth}{0.5pt}
\renewcommand{\footrulewidth}{0pt}
\AtBeginDocument{%
  \renewcommand{\headrule}{{\color{colRule}\hrule width\headwidth height 0.5pt}}%
}

\setlist[itemize]{leftmargin=1.2em, itemsep=1pt, topsep=2pt, parsep=0pt}
\setlist[enumerate]{leftmargin=1.2em, itemsep=3pt, topsep=2pt, parsep=0pt}

\setlength{\parskip}{2pt}
\setlength{\columnsep}{10pt}

\def\CAPETBEGINFILE{\begin{document}}
\def\CAPETENDFILE{\end{document}}
\else
\def\CAPETBEGINFILE{}
\def\CAPETENDFILE{}
\fi

\ficheHeaders{T24}{Systèmes embarqués}{colElec}

\CAPETBEGINFILE

% -----------------------------------------------------------------------
%  BLOC EN-TETE
% -----------------------------------------------------------------------

\ficheEntete{T24}{Systèmes embarqués - Microcontrôleurs, API, FPGA}{colElec}%
  {électronique - Commande numérique et systèmes programmables}%
  {\textbf{Présence :} \textbf{6/13} sujets (2013-2026) - thème \emph{fréquent}}%
  {microcontrôleur (Arduino, PIC) - registre ADCON0 - API S7-1500 - PWM - bus CAN/OneWire/LoRa - DSP - FPGA/VHDL (programme) - mise à l'échelle automate}

% -----------------------------------------------------------------------
%  1. RAPPELS DE COURS
% -----------------------------------------------------------------------

\begin{tcolorbox}[
  ficheSection=colElec,
  title={1. Rappels de cours},
]

% -- 1.1 Architecture systemes embarques -------------------------------
\textbf{1.1 Architecture des systèmes embarqués}

\smallskip
\begin{center}
\renewcommand{\arraystretch}{1.25}
\begin{tabular}{llll}
\toprule
\textbf{Composant} & \textbf{Rôle} & \textbf{Exemples CAPET} & \textbf{Prog.} \\
\midrule
Microcontrôleur & CPU + périphériques intégrés & Arduino, PIC 18F8723 & C/Python \\
DSP & Traitement signal haute vitesse & DSP 40 MHz (2024) & C optimisé \\
API/PLC & Contrôle industriel robuste & Siemens S7-1500 (2024) & Ladder/Grafcet \\
FPGA & Circuit reconfigurable parallèle & (Programme SII) & VHDL \\
\bottomrule
\end{tabular}
\end{center}

\medskip\hrule\vspace{4pt}

% -- 1.2 Microcontroleur - CAN interne --------------------------------
\textbf{1.2 Microcontrôleur - CAN interne (PIC 18F8723)}

\begin{itemize}
  \item 1 CAN partagé entre 16 entrées analogiques AN0→AN15 via multiplexeur
  \item Sélection de la voie : registre ADCON0, champs CHS3:CHS0
  \item Résolution 10 bits ; Quantum $q = V_{ref}/2^{10} = V_{ref}/1024$
  \item Démarrage conversion : bit GO/DONE = 1
\end{itemize}

\medskip\hrule\vspace{4pt}

% -- 1.3 PWM ----------------------------------------------------------
\textbf{1.3 PWM - Modulation de Largeur d'Impulsion}

Signal créneau à fréquence fixe (50 Hz pour servomoteur), rapport cyclique $\alpha = t_{on}/T$.

\smallskip
\begin{center}
\begin{tikzpicture}[scale=0.85, font=\footnotesize]
  \draw[->] (0,0) -- (6.3,0) node[right] {$t$};
  \draw[->] (0,-0.3) -- (0,2.0) node[above] {$V_{PWM}$};
  \draw[thick, colElec] (0,1.4) -- (1.5,1.4) -- (1.5,0) -- (6,0);
  \draw[dotted] (1.5,0) -- (1.5,1.4);
  \node[below] at (0.75,0) {\scriptsize $t_{on}$};
  \draw[dashed] (6,0) -- (6,-0.3);
  \node[below] at (6,-0.3) {\scriptsize $T = 20\,\mathrm{ms}$ (50 Hz)};
\end{tikzpicture}
\end{center}

\[
  t_{on} = \frac{\theta - \theta_{min}}{\theta_{max} - \theta_{min}}\,(t_{on,max} - t_{on,min}) + t_{on,min}
\]

Fréquence d'échantillonnage pour résolution angulaire $\delta\theta$ :
\[
  f_{ech} = \frac{\theta_{max} - \theta_{min}}{\delta\theta \cdot T} = \frac{N_{positions}}{T}
  \quad \Rightarrow \quad
  f_{ech} = N_{pos} \times f_{PWM}
\]
Exemple (2025) : $\delta\theta = 0{,}1^\circ$, plage $180^\circ$ - $N = 1800$ - $f_{ech} = 90\,\mathrm{kHz}$ (PWM à 50 Hz).

\medskip\hrule\vspace{4pt}

% -- 1.4 API Siemens S7-1500 ------------------------------------------
\textbf{1.4 Automate S7-1500 - Mise à l'échelle}

\[
  T = T_{min} + \frac{PIW8 - 0}{27\,648 - 0}\,(T_{max} - T_{min})
\]
\begin{itemize}
  \item Module AI : PIW8 $\in [0\,;\,27\,648]$ pour $I_s \in [4\,;\,20]\,\mathrm{mA}$
  \item Programme cyclique : lecture entrées - traitement - écriture sorties
  \item Capteur OneWire DS1820 : commande 0x44 (conversion) ; 0xBE (lecture scratchpad)
\end{itemize}

\medskip\hrule\vspace{4pt}

% -- 1.5 DSP ----------------------------------------------------------
\textbf{1.5 DSP - Calcul de fréquence maximale}

\[
  f_{max} = \frac{P \cdot n}{60}
  \quad P : \text{points/tour codeur}, \quad n : \text{vitesse (tr/min)}
\]
Condition : $f_{DSP} \geq f_{max}$. Exemple (2024) : $f_{max} = 1326\,\mathrm{Hz}$, $f_{DSP} = 40\,\mathrm{MHz}$.

\medskip\hrule\vspace{4pt}

% -- 1.6 Bus CAN Controller Area Network -------------------------------
\textbf{1.6 Bus CAN (Controller Area Network)}

Bus série différentiel multi-maîtres, débit jusqu'à 1 Mbit/s.\\
Structure de trame : SOF (1) + ID (11) + RTR (1) + IDE+r0 (2) + DLC (4) + Data (0-64) + CRC (15) + délimiteur CRC (1) + ACK (2) + EOF (7) = 108 bits max (hors bits de stuffing).

\end{tcolorbox}

% -----------------------------------------------------------------------
%  2. FORMULAIRE SYNTHETIQUE
% -----------------------------------------------------------------------

\begin{tcolorbox}[
  ficheSection=colElec,
  title={2. Formulaire synthétique},
]
\small
\begin{multicols}{2}
\renewcommand{\arraystretch}{1.45}
\begin{tabular}{@{}p{5cm}l@{}}
\toprule
\textbf{Formule} & \textbf{Usage} \\
\midrule
\multicolumn{2}{@{}l}{\footnotesize\textit{PWM}} \\
$\alpha = t_{on}/T$ & Rapport cyclique \\
$f_{ech} = N_{pos} \times f_{PWM}$ & Résolution angulaire \\
$N_{pos} = (\theta_{max}-\theta_{min})/\delta\theta$ & Nb positions \\
\midrule
\multicolumn{2}{@{}l}{\footnotesize\textit{DSP}} \\
$f_{max} = P \cdot n / 60$ & Fréquence codeur \\
Condition : $f_{DSP} \geq f_{max}$ & Choix DSP \\
\bottomrule
\end{tabular}

\columnbreak

\renewcommand{\arraystretch}{1.45}
\begin{tabular}{@{}p{5cm}l@{}}
\toprule
\textbf{Formule} & \textbf{Usage} \\
\midrule
\multicolumn{2}{@{}l}{\footnotesize\textit{API S7-1500}} \\
$T = T_{min} + PIW8/27648 \cdot \Delta T$ & Mise à l'échelle \\
PIW8 $\in [0\,;\,27\,648]$ & Plage module AI \\
\midrule
\multicolumn{2}{@{}l}{\footnotesize\textit{Bus CAN}} \\
Débit max : 1\,Mbit/s & CAN 2.0 \\
Trame max : 108\,bits (8 octets data) & Structure \\
$D_b = 1/T_b$ & Débit depuis chrono. \\
\bottomrule
\end{tabular}
\end{multicols}
\normalsize
\end{tcolorbox}

% -----------------------------------------------------------------------
%  3. QUESTIONS TYPES
% -----------------------------------------------------------------------

\begin{tcolorbox}[
  ficheSection=colElec,
  title={3. Questions types - extraites des sujets réels},
]

% -- QT1 -------------------------------------------------------------

\begin{tcolorbox}[questiontype,
  title={QT1 - Registre ADCON0 du PIC - configuration CAN \hfill\normalfont\footnotesize E1/E2 2022 (PIC 18F8723)}
]
\textbf{Énoncé.}
Pour sélectionner l'entrée AN2 du PIC 18F8723 et lancer la conversion, remplir le registre ADCON0 champ par champ. Donner la valeur hexadécimale.

\tcbline

\textbf{Corrigé.}

\textbf{Étape 1 — Structure ADCON0.} [ADON | GO/DONE | CHS3 | CHS2 | CHS1 | CHS0 | - | ]
\begin{itemize}
  \item ADON = 1 (CAN activé)
  \item GO/DONE = 1 (lancer la conversion)
  \item CHS3:CHS0 = 0010 (voie AN2)
\end{itemize}

\textbf{Étape 2 — Assemblage.} Registre : $1\,1\,0010\,00_2$ (si les 2 LSB non utilisés = 0) = $\mathrm{C8_{16}}$ (à vérifier sur le datasheet exact).

\textbf{Conclusion.} Le codage binaire et la méthode de configuration bit par bit sont corrects ; la valeur hexadécimale finale reste \textbf{[corrigé reconstitué]} tant que la table officielle ADCON0 fournie dans le sujet n'est pas confirmée — l'attitude attendue est de le signaler plutôt que d'affirmer un résultat non vérifiable.
\end{tcolorbox}

\vspace{4pt}

% -- QT2 -------------------------------------------------------------

\begin{tcolorbox}[questiontype,
  title={QT2 - PWM servomoteur - fréquence d'échantillonnage \hfill\normalfont\footnotesize E2 Commune 2025 (Q18-Q19)}
]
\textbf{Énoncé.}
Signal PWM à 50 Hz pour un servomoteur ($-90^\circ$ à $+90^\circ$). Précision angulaire requise : $\delta\theta = 0{,}1^\circ$. À quelle fréquence le récepteur doit-il échantillonner le signal PWM ? Le microcontrôleur Arduino (84 MHz) est-il compatible ?

\tcbline

\textbf{Corrigé.}

\textbf{Étape 1 — Nombre de positions.} $N = \dfrac{180^\circ}{0{,}1^\circ} = 1800$.

\textbf{Étape 2 — Fréquence d'échantillonnage nécessaire.}
\[
  f_{ech} = N \times f_{PWM} = 1800 \times 50 = \mathbf{90\,000\,\mathrm{Hz}} = 90\,\mathrm{kHz}
\]

\textbf{Étape 3 — Vérification Arduino.} $f_{horloge} = 84\,\mathrm{MHz} \gg 90\,\mathrm{kHz}$.

\textbf{Conclusion.} L'exigence d'échantillonnage est très largement respectée par l'horloge Arduino : la solution est \textbf{conforme au cahier des charges} sur ce critère, à condition que le code de mesure du temps haut reste suffisamment optimisé pour ne pas introduire de latence logicielle.
\end{tcolorbox}

\vspace{4pt}

% -- QT3 -------------------------------------------------------------

\begin{tcolorbox}[questiontype,
  title={QT3 - Mise à l'échelle automate - programme de conversion \hfill\normalfont\footnotesize E1 2024 (S7-1500)}
]
\textbf{Énoncé.}
Le module AI du S7-1500 reçoit un signal 4-20\,mA d'un thermocouple type J ($T \in [-20\,;\,+200]\,^\circ\mathrm{C}$). PIW8 $\in [0\,;\,27\,648]$. Écrire la relation de mise à l'échelle.

\tcbline

\textbf{Corrigé.}

\textbf{Étape 1 — Relation affine.}
\[
  T = T_{min} + \frac{PIW8 - 0}{27\,648 - 0}\,(T_{max} - T_{min})
  = -20 + \frac{PIW8}{27\,648} \times 220
\]

\textbf{Étape 2 — Implémentation.} En Ladder/Blocs fonctionnels S7-1500 : instruction SCALE ou calcul avec blocs ADD, MUL, DIV avec les coefficients $a = 220/27648$ et $b = -20$.

\textbf{Étape 3 — Vérification aux bornes.} $PIW8 = 0 \Rightarrow T = -20^\circ$C ; $PIW8 = 27\,648 \Rightarrow T = 200^\circ$C.

\textbf{Conclusion.} La loi affine est validée aux deux bornes de la plage : la conversion respecte le cahier des charges de mesure sur l'intégralité de l'étendue utile du capteur.
\end{tcolorbox}

\end{tcolorbox}

% -----------------------------------------------------------------------
%  4. ATTENTES DU JURY
% -----------------------------------------------------------------------

\begin{tcolorbox}[
  jury,
  title={4. Attentes du jury - Erreurs récurrentes et points de vigilance},
]

\begin{enumerate}[label=\textbf{\arabic*.}, leftmargin=*, itemsep=5pt]

  \item \textbf{Configurer un registre bit par bit.} Justifier chaque bit de ADCON0 (ou tout autre registre de microcontrôleur) - écrire la valeur hex sans justification est insuffisant.

  \item \textbf{Câblage d'une entrée/sortie automate.} Jury 2021 : "Un candidat au CAPET doit être capable de câbler une entrée TOR/analogique d'un API industriel."

  \item \textbf{Confusion API vs microcontrôleur.} L'API est robuste, programmé en Ladder/Grafcet pour l'industrie. Le microcontrôleur est compact, programmé en C/Python pour les fonctions embarquées légères. Ne pas interchanger.

  \item \textbf{Mise à l'échelle API.} Ne pas oublier le décalage (offset) quand la plage ne commence pas à zéro (ex. $PIW8 = 0 \Rightarrow T = -20^\circ$C). Vérifier les deux bornes.

  \item \textbf{Bus CAN vs CAN (convertisseur).} Dans les sujets, les deux coexistent (2026 E1). Distinguer par le contexte : bus de terrain vs convertisseur analogique-numérique.

  \item \textbf{FPGA/VHDL.} Absent des sujets réels 2013-2026 mais au programme. Si la question apparaît : entité + architecture, signal, process concurrent vs séquentiel.

\end{enumerate}

\medskip
\begin{tcolorbox}[
  colback=colJuryFrame!12!colJury,
  colframe=colJuryFrame,
  left=4pt, right=4pt, top=2pt, bottom=2pt,
  fontupper=\small\itshape
]
Rapport jury 2023 : «~Notions d'algorithmes, variables informatiques, boucles, systèmes embarqués sont des connaissances de base attendues des candidats.~»\\
Rapport jury 2021 : «~Un candidat au CAPET doit être capable de câbler une entrée ou une sortie TOR/analogique d'un API industriel programmable.~»
\end{tcolorbox}

\end{tcolorbox}

% -----------------------------------------------------------------------
%  PIED - PRESENCE DANS LES SUJETS
% -----------------------------------------------------------------------

\begin{tcolorbox}[
  ficheSection=colPresence,
  title={Présence dans les sujets analysés (2013-2026)},
]
\begin{center}
\footnotesize\renewcommand{\arraystretch}{1.15}
\begin{tabular}{lccccccccccccc}
\toprule
\textbf{Année} & 13 & 14 & 15 & 16 & 17 & 18 & 19 & 20 & 21 & 22 & 23 & 24 & 25 \\
\midrule
E1 (commune)   & - & $\checkmark$ & - & - & - & - & - & - & $\checkmark$ & $\checkmark$ & - & $\checkmark$ & - \\
E2 discipl. IE & - & - & - & - & - & - & $\checkmark$ & - & - & - & - & $\checkmark$ & $\checkmark$ \\
E2 appl. comm. & - & - & - & - & - & - & - & - & - & - & - & $\checkmark$ & $\checkmark$ \\
\bottomrule
\end{tabular}
\end{center}
\smallskip
\footnotesize\textbf{Sous-thèmes fréquents :}
registre CAN microcontrôleur (PIC ADCON0) - OneWire DS1820 - mise à l'échelle API S7-1500 - PWM servomoteur - DSP fréquence - bus CAN trame - LoRaWAN IoT

\medskip
\textbf{Ordres de grandeur (repères rapides).}
\begin{itemize}[leftmargin=1.3em]
  \item Réseau BT FR : tension phase-neutre typique \SI{230}{\volt}, tension phase-phase typique \SI{400}{\volt}, fréquence \SI{50}{\hertz}.
  \item Courants d'étude fréquents en sujet : de quelques \si{\ampere} à quelques dizaines de \si{\ampere}.
  \item En conversion d'énergie, un rendement attendu reste généralement entre 0{,}8 et 0{,}98 selon la technologie.
\end{itemize}

\medskip
\textbf{Synthèse express (à refaire sans support).}
\begin{enumerate}[leftmargin=1.4em]
  \item Poser les hypothèses du modèle et les conventions de signe.
  \item Écrire 2 à 4 relations clés avec unités.
  \item Réaliser un calcul numérique complet et vérifier l'ordre de grandeur.
  \item Justifier le résultat par un argument technique de niveau CAPET.
\end{enumerate}

\end{tcolorbox}

\CAPETENDFILE
